\section{Background}\label{sec:background}

Large neural networks have produced state-of-the-art results in numerous fields such as computer vision
 or natural language processing. To deploy these neural network model in an industry application, the most common option is training 
 the model in a data center utilizing huge GPU clusters, liquid cooling, etc. At inference time, the clients could also access the model in the data center 
 via internet.

 However, with the proliferation of Internet-of-Things (IoT) devices and applications, neural network models are required to run more efficiently on
 resource-constrained devices to facilitate the growing number of edge AI use cases \cite{edgeai}. Efficiency or compactness refers to the model giving outputs accurate
 enough to enable the specific application, while also inducing low memory and energy-consumption overhead. The latter two translates to the model having
 a small parameter count and an architecture which minimizes latency and floating-point operations (FLOPs). Depending on the deployment scenario, a model may
 receive inputs that are far from the distribution of the training data or inputs that are manipulated to make the model produce a wrong output. Robustness of 
 the model refers to its ability to maintain its performance even on such inputs.  To sum up, a trained neural network has to satisfy accuracy, compactness and robustness requirements in order to become qualified for industrial
  applications. 

  In academia, there has been extensive research to improve them, but most works tackled them separately \cite{effnet} \cite{pruning_survey24}.
  To improve accuracy, the general approach was training overparameterized networks with regularization. Model compression techniques (i.e. weight pruning, knowledge
  distillation) and novel architecture have been proposed to achieve compactness while retaining accuracy. Finally, increasing robustness has been achieved by 
  augmenting the training dataset to better represent the distribution of inputs at inference \cite{effectofadvtraining}.

  However, meeting these requirements simultaneously - so not increasing one at the expense of another - was found to be a challenge \cite{cards}.
  Therefore, it is important to evaluate a technique tackling one of the three requirements on all three aspects. In my case, measuring the impact of model
  compression on accuracy and robustness as well. In the following sections, we give an overview of the techniques relevant to this work. 

\subsection{Global Magnitude-based Pruning}\label{subsec:prune}

There are many contemporary methods for model compression. The main four categories include network pruning, knowledge distillation, tensor decomposition,
and quantization \cite{survey}, although it is also common to propose hybrid strategies. Network pruning refers to evaluating weights or groups of weights
of the model based on a certain importance criteria and removing the ones which were considered the least important. This importance signal could weight
magnitude, the sensitivity of loss with respect to the weight, or other heuristics. After pruning, a fine-tuning or retraining phase is applied using the 
remaining weights in order to refit the pruned model to the training data.
When the magnitude is utilized for representing the importance of weights, the pruning step uses a threshold for comparison, and every weight with magnitude
below the threshold gets removed. Global pruning refers to having only one threshold and using it across all layers, regardless of their shape and function
(i.e. classifier layer, convolutional layer). If every layer has its separate, local threshold, the pruning is applied uniformly across the entire network,
which can lead to sub-optimal sparsity and the increased complexity of tuning these thresholds for a deep network \cite{pruning_survey24}. In contrast, 
global pruning enables layers with different roles (i.e. feature extractor, classifier) to contribute to the overall sparsity non-uniformly.

\subsection{Weight and Learning Rate Rewinding}\label{subsec:rewind}

\subsection{Sparsity-enhancing Regularization}
In the field of deep learning, regularization aims at reducing the generalization error of the model, often at the expense of the training error
\cite{Goodfellow}. This mostly involves some modification made to the training algorithm, the training dataset or the way the output model is used. 
One common regularization technique is adding a norm penalty term to the objective function which expresses a prior belief about the weights of the
final network. Injecting some prior knowledge into the training process could also be employed via adding constraints to the optimization algorithm, 
augmenting the dataset by applying transformations on the existing samples which the model must be invariant to, or sharing weights across different
parts of the network.    

As the use of norm penalties for model compression purposes is an established practice, I used the same approach. The general form of the cost
function $C$ in the case of supervised learning can be written as: 
\begin{equation}
  C(W) = \sum_{i} \mathcal{L}(f(x_i, W), d_i) + \lambda \cdot \mathcal{H}(W)
\label{eq:reg_cost}
\end{equation}

where $W$ denotes the parameters of the neural network, $\mathcal{L}$ denotes the loss function used for training, function $f$ is the output of the
network, $x_i$ is the training sample and $d_i$ is the corresponding target. The regularization term consists of a function of the weights denoted 
by $\mathcal{H}(W)$, and a coefficient $\lambda$ which determines the strength or importance of regularization compared to the performance of the 
model on the training samples.
 
In case of network pruning, our general assumption is that there is a sparse subnetwork inside the dense neural network which could achieve the same 
accuracy as the dense network when retrained in isolation (which means training it further after dropping the rest of the units from the dense 
network) \cite{winningticket}. Hence, the penalty term must represent the belief that the number of nonzero weights in the initial dense network 
should be as high as possible. Although finding the aforementioned sparse subnetwork is the primary goal of network pruning, the penalty term 
should not dominate the loss term in the objective to enable learning the task at hand.

\subsection{Iteratively Reweighted $l_1$ Minimization}\label{subsec:irl1}

As L0 regularization results in a non-differentiable cost function, it is common to use convex relaxtions of it, such as L1 regularization.
Besides L1, there are non-convex functions which better approximate the L0 term, such as the log-sum function visualized in Figure \ref{fig:logsum}.

\begin{figure}[H]
    \includegraphics[width=0.5\linewidth]{figs/logsum.png}
    \centering
    \caption{Comparison of L0, L1 and log-sum functions.}
    \label{fig:logsum}
\end{figure}

As the figure shows, the \textit{log-sum} function approximates the L0 penalty term better; it penalizes smaller weights more,
 as the magnitude of the gradient increases as it approaches zero (from both sides).
Boyd et. al. \cite{boyd} proposes an iterative algorithm for reweighted L1 minimization, which will be referred to as 
\textit{Iterative Reweighted $l_1$ Minimization } (IRL1) in this document. The subject of this paper is Compressed Sensing which is a subfield 
of signal processing aiming to reconstruct a signal vector $x_0 \in R^n$ from measurements $y = \Phi x_0$ where $\Phi \in R^{m\times n}$
and $m < n$. As there are fewer measurements than the dimension of the unknown signal, there are infinite number of possible solutions to $x_0$.
The main result of compressed sensing is that the signal $x_0$ can be reconstructed if it is sparse in some domain. There are more conditions
that have to be met to reconstruct the sparse signal, but these go beyond the scope of this work. The reconstruction is done by minimizing the
L1 norm of $x_0$ while using equations of $y=\Phi x_0$ as constraints. Authors aimed to improve upon the utilization of the L1 norm in terms of the sparsity
achieved in the reconstructed signal. They showed that solving a sequence of weighted L1 minimization problems finds a local minimum of the concave log-sum
function. As previously mentioned, the log-sum function is a better approximation of the L0 term, so it is more effective in sparsifying the
reconstruction than the L1 term. Thus, the adoption of this algorithm on the regularization term used for network pruning could potentionally find a
smaller subnetwork in the initial dense network.  

As the experiments will be conducted on a convolutional network, the use of IRL1 is described in terms of weight tensors of a convolutional layer.
The weighted L1 term can be written as:   
\begin{equation}
 \mathcal{H}(W^{(\ell)}) = \sum_{l,i, j} \sum_{u,v} A^{(\ell)}_{i,j,u,v} \odot\big| W^{(\ell)}_{i,j}(u,v) \big|
\label{eq:wl1_term}
\end{equation}

where $A^{(\ell)}$ is the tensor of weights used to balance the penalty induced by network weight magnitudes in layer $\ell$. 
It has the same dimension as the weight tensor of the convolutional layer and $\odot$ denotes the element-wise multiplication between tensors
$A^{(\ell)}$ and $W^{(\ell)}$.  As both the weights of the network and the weights of the regularization get updated, let $k$ denote the current
iteration. Before the first iteration, all regularization weights are initialized as 1. Then after a few gradient descent steps, these weights
are updated as follows:
\begin{equation}
 A_{i,j,u,v}^{(\ell, k+1)} = \frac{1}{\big|W^{(\ell, k)}_{i,j,u,v} \big| + \epsilon}
\label{eq:weight_update}
\end{equation}

where $\epsilon$ represents a small nonzero number which prevents division by zero and balances the extent to which small weights are penalized. 

Figure \ref{fig:epsilon} depicts three final weight distributions of a convolutional layer corresponding to IRL1 runs with different 
$\epsilon$ values. The additional distribution at the left shows the final distribution when training the layer without any regularization.

\begin{figure}[H]
    \includegraphics[width=1\linewidth]{figs/eps2.png}
    \centering
    \caption{Final weight distributions of a convolutional layer corresponding to IRL1 runs with different $\epsilon$ values.}
    \label{fig:epsilon}
\end{figure}

Not surprisingly, the distribution in blue resembles a normal distribution. IRL1 regularization with the largest value of $\epsilon$
pushed the whole distribution towards zero, similarly to L1 regularization. On the other hand, the rightmost distribution corresponding to the 
smallest $\epsilon$ value has produces a spike in a small zero-centered interval, thus resembling the application of a hard thresholding operator.
The hard thresholding operator is usually how L0 regularization is implemented in the context of network pruning. 
Hence, the greater the value of $\epsilon$, the better the IRL1 algorithm resembles L1 minimization, while smaller $\epsilon$ values make the
effect of IRL1 more similar to L0 minimization. 

The experiments done in \cite{boyd} kept $\epsilon$ constant, but dynamically adjusting it to account for the dynamics of training could also be a valid
strategy. 

Another hyperparameter of this method is the frequency of updating tensor $A$. There are no guidelines for it in \cite{boyd}, as it is done
after convergence of the weighted L1 minimization. In our case, the non-convex objective function of training the neural network never converges,
so an update frequency has to be set heuristically. If the L1 term is reweighted too often, it may destabilize training as the optimizer would 
face a frequently-changing objective function. On the other hand, if the update interval is too long, the optimization may find a dense solution 
with training error too low for the reweighted L1 term to induce sparsity. Different sparsity curves - captured during the proposed pruning process in
 Chapter \ref{sec:method} -
corresponding to different update frequencies are compared in Appendix \ref{sec:appendix} to provide some intuition on the effect of this hyperparameter.
